% !Mode:: "TeX:UTF-8"
%%%论文封面
%%  教育硕士专业学位是按   华南师范大学教师教育学部 2024 年 6 月 6 日颁布的华南师范大学研究生学位论文撰写规范  所制定
%%  学术硕士专业学位是按 华南师范大学数学科学学院 2024 年 5 月 23 日颁布年颁布的 研究生学位论文撰写规范（适用学科：0701数学） 所制定

%%%%%%%%%%%%%%%%%%%%%%%%%%%%%%
% 专业学位     中文封面内容
%% “论文分类号”按《中国图书资料分类法》的分类号填写
%%% 具体可查询 http://ztflh.xhma.com
% 分类号 学术型:
% 计算数学 O24   应用数学 O29
% 概率论与数理统计 O21   运筹学  O22
% 统计学 C8   课程与教学论  G427

\classification{}  % 分类号
\UDC{} % UDC: 表示国际十进分类法(Universal Decimal Classification),是世界一般报告或者论文的封面在左上角注明分类号，便于信息交换和处理。
%%%  专业学位的 分类号、UDC：不填写
%%%  学术学位的  分类号: 有  计算数学 O24   应用数学 O29   概率论与数理统计 O21    运筹学  O22 统计学 C8   课程与教学论  G427
\confidential{公开}
%% “密级：学位论文密级分为公开、内部、涉密。内部论文和涉密论文严格按照学校审批的保密级别和保密期限填写，如内部*年、秘密*年；其他论文均填公开
\studentID{20222005253}
%%   填写学位申请人的学号

\type{本~科}

%%% ====================
%%% ====== 中文标题 ======
%%% ====================
%%% 论文标题: 应准确概括整篇论文的核心内容，简明扼要，一般不超过25个汉字。必要时可加副标题。
\title{\CASTunderline[\textwidth]{TinyTheorem：}
  \\ \vspace{0.18cm}
\CASTunderline[\textwidth]{一个依值类型的编程语言暨定理证明器原型}}
%% 如上是封面中的 论文中文标题, 需要人工分行

\titleab{基于依赖类型论的编程语言暨定理证明器实现}
%\subtitle{}    % 中文副标题,  排版时自动居右, 这是可选项, 若不需要可注解掉

%%% ====================
%%% ====== 英文标题 ======
%%% ====================
%%% 英文标题

\etitle{\CASTunderline[\textwidth]{TinyTheorem: }
  \\ \vspace{0.18cm}
  \CASTunderline[\textwidth]{A prototype for a dependently-typed}
  \\ \vspace{0.18cm}
  \CASTunderline[\textwidth]{programming language with theorem proving}
}

%%%==================== ==========
%%%==========  学位申请人 ==========
\author{高振明}
\eauthor{GAO Zhenming}
%%% 英（外）文封面作者姓名按姓在前、名在后的格式书写，姓名需写全拼，如LI Jianguo

%%%==================== ==========
%%%==========  学科专业  ==========
%%% 学科专业: 申请人所在学科专业，必须严格按参照学科代码/专业学位代码字典中的二级学科名称/专业学位领域名称填写，若所选一级学科/专业学位类别下无二级学科/专业学位领域，则填写一级学科/专业学位类别名称，不可用简称。

%%%    学术学位
%  \major{计算数学}
% \emajor{Computional Mathematics}
%  \major{应用数学}
% \emajor{Applied Mathematics}
%  \major{基础数学}
% \emajor{Pure Mathematics}
%  \major{概率论与数理统计}
% \emajor{Probability Theory and Mathematic Statistics}
%  \major{运筹学与控制论}
% \emajor{Operations Research and Cybernetics}

%%%       专业
\major{软件工程}
\emajor{Software Engineering}

%%%==================== ==========
%%%==========   所在学院  ==========
%%% 申请人所在二级培养单位的名称，必须填写与学校官网一致的规范名称，不可用简称。
\college{人工智能学院}
\ecollege{School of Artificial Intelligence}

%%%==================== ==========
%%%==========  导师姓名及职称 ==========
%%%  填写指导教师的姓名与职称。若有两名指导教师，第一（校内）导师姓名和职称写在前面，第二（或校外）导师姓名和职称写在后面，中间空一格。第一（校内）导师必须与研究生管理系统一致。
\advisor{邹竞辉} % 导师姓名
\ProfessionalTitle{讲师} % 及职称
\eadvisor{ZOU Jinhui}  %

%%%==================== ==========
%%%==========  论文完成时间 ==========
% 填写学位论文最终定稿打印成文的年月，封面的日期用汉字书写，如二〇二四年六月、二〇二四年十二月
\submityear{二〇二五} % 年, 用于中文封面
\submitmonth{五}         % 月, 用于中文封面
\submitengyear{2025}  % 年, 用于英文封面
\submitengmonth{MAY} % 月, 用于英文封面
\gradyear{2025}
\gradmonth{7}

\maketitle   % 生成中文封面

\strut\newpage

\begin{cabstract}
  （..略..）
\end{cabstract}

%% 中文关键词, 各关键词之间用逗号隔开, 最后一个关键词不加标点.
%% 关键词一般在3—8个之间.
\begin{ckeywords}
  逻辑推理素养, 高考评价体系, 圆锥曲线, 培养策略
\end{ckeywords}

\begin{center}
  \vskip \stretch{1}
  {
    \begin{spacing}{1.5}
      \songti\bf\erhao {TinyTheorem：}
      \\ \vspace{0.18cm}
      一个依值类型的编程语言暨定理证明器原型
  \end{spacing} } % 论文题目
\end{center}
\vspace{0.68cm}

\makecabstract

\cleardoublepage

\begin{center}
  \vskip \stretch{1}
  {
    \begin{spacing}{1.68}
      \songti\bfseries\erhao
      TinyTheorem:
      \\ \vspace{0.18cm}
      A prototype for a dependently-typed
      \\ \vspace{0.18cm}
      programming language with theorem proving
  \end{spacing} } % 论文题目
  \vspace{0.68cm}
  %\vskip \stretch{1}
\end{center}


\begin{eabstract}
  (..omitted..)
\end{eabstract}

%% 英文关键词, 各关键词之间用逗号隔开. 外文关键词应与中文关键词相对应
\begin{ekeywords}
  logical reasoning literacy, college entrance examination evaluation system, conic curve, cultivation strategies.
\end{ekeywords}
\makeeabstract